\begin{abstract}

Artifacts in electroencephalography (EEG) hinder automated analysis, particularly in seizure detection. While EEG data are costly to label at scale, modern generative models can synthesise realistic, label-aware artifact segments suitable for augmentation. A latent diffusion-based augmentation method with a per-class training strategy is proposed to improve artifact classification performance in imbalanced settings. To address class imbalance, the training data is extended with synthetic EEG signal segments, improving data diversity and class representation. Using the TUH EEG Artifact (TUAR) corpus, the proposed method is evaluated and compared against existing non-generative and generative augmentation techniques. Results show that latent diffusion-based augmentation improves artifact classification by 3.7\% in Macro F1 and 3.5\% in Macro PR AUC and demonstrates small but consistent improvements when combined with existing augmentation pipelines. Although the proposed method is not able to outperform all baseline approaches, it demonstrates consistent improvements in classification performance and showcases the potential of latent diffusion-based augmentation for imbalanced EEG artifact datasets. 

\end{abstract}